% Paper I: Fejér–heat generators and Lipschitz control for the Weil quadratic functional
% Style: Rodgers–Tao (Forum of Mathematics, Pi)
\documentclass[11pt,a4paper]{article}

\usepackage{amsmath,amssymb,amsthm}
\usepackage{mathrsfs}
\usepackage{hyperref}
\usepackage{enumitem}
\usepackage{booktabs}
\usepackage{geometry}
\geometry{margin=1in}

% Theorem environments (Tao style)
\newtheorem{theorem}{Theorem}[section]
\newtheorem{lemma}[theorem]{Lemma}
\newtheorem{proposition}[theorem]{Proposition}
\newtheorem{corollary}[theorem]{Corollary}
\theoremstyle{definition}
\newtheorem{definition}[theorem]{Definition}
\theoremstyle{remark}
\newtheorem{remark}[theorem]{Remark}

% Custom commands
\newcommand{\R}{\mathbb{R}}
\newcommand{\C}{\mathbb{C}}
\newcommand{\Z}{\mathbb{Z}}
\newcommand{\T}{\mathbb{T}}
\newcommand{\N}{\mathbb{N}}
\DeclareMathOperator{\supp}{supp}

\title{Fej\'er--heat generators and Lipschitz control\\for the Weil quadratic functional}

\author{
  {\sc First Author}\thanks{Department of Mathematics, University A; email: author1@university.edu}
  \and
  {\sc Second Author}\thanks{Department of Mathematics, University B; email: author2@university.edu}
}

\date{}

\begin{document}

\maketitle

\begin{abstract}
The Guinand--Weil explicit formula expresses a quadratic functional $Q$ on test function spaces via two competing terms: an Archimedean integral encoding digamma contributions and a prime sum over von Mangoldt weights at logarithmically-spaced nodes. This paper establishes foundational properties of $Q$ that are prerequisites for spectral positivity arguments. We first fix a canonical normalization (T0) that matches the classical Guinand--Weil formulation under the change of variables $\eta = 2\pi\xi$. We then introduce Fej\'er--heat generators---products of Fej\'er triangular windows with Gaussian heat kernels---and prove that their finite non-negative combinations form a dense cone in the space of even, non-negative test functions on any compact frequency window $[-K, K]$ (Theorem A1'). The density argument proceeds via a three-step approximation: mollification by the heat kernel, positive Riemann sum decomposition, and Fej\'er truncation with controlled error. Finally, we establish Lipschitz continuity of $Q$ on compact windows (Lemma A2), showing that the functional is continuous in the uniform norm with explicit Lipschitz constant depending on the Archimedean density and the finite prime weight sum. The local finiteness of the prime sampler---only finitely many prime nodes fall within any bounded interval---is central to this continuity. These results, together with the spectral analysis of the Toeplitz barrier and RKHS prime cap developed in [Part~II], form the analytic foundation for studying positivity properties of the Weil functional. This paper is the first of three developing an operator-theoretic approach to the Weil positivity criterion.
\end{abstract}

\tableofcontents

%=============================================================================
\section{Introduction}
%=============================================================================

\subsection{The Weil functional}

The Guinand--Weil explicit formula~\cite{Guinand1948,Weil1952} provides a remarkable connection between the prime distribution and the behavior of $L$-functions through a quadratic functional on test function spaces. For a suitable even test function $\Phi : \R \to \R$, the \emph{Weil quadratic functional} takes the form
\begin{equation}\label{eq:Q-def}
  Q(\Phi) := \int_{\R} a_*(\xi)\, \Phi(\xi)\, d\xi
           - \sum_{n \geq 2} w_Q(n)\, \Phi(\xi_n),
\end{equation}
where:
\begin{itemize}
\item The \emph{Archimedean density} is $a_*(\xi) = 2\pi\bigl(\log\pi - \Re\psi(\tfrac{1}{4} + i\pi\xi)\bigr)$, with $\psi$ the digamma function.
\item The \emph{prime nodes} are $\xi_n = \frac{\log n}{2\pi}$ for $n \geq 2$.
\item The \emph{Weil weights} are $w_Q(n) = \frac{2\Lambda(n)}{\sqrt{n}}$, where $\Lambda$ is the von Mangoldt function.
\end{itemize}

The functional $Q$ encodes deep information about the distribution of prime numbers. The Archimedean term is a smooth integral against a known density, while the prime term is a discrete sum sampling the test function at logarithmically-spaced nodes. Understanding when $Q(\Phi) \geq 0$ for appropriate test classes is a fundamental problem in analytic number theory.

\subsection{Main results}

The goal of this paper is to establish three foundational properties of $Q$:

\begin{enumerate}[label=(\alph*)]
\item \textbf{Normalization (T0):} A canonical correspondence between our frequency coordinate $\xi$ and the classical Guinand--Weil variable $\eta = 2\pi\xi$.

\item \textbf{Cone density (A1'):} The Fej\'er--heat generators form a dense cone in the space of even, non-negative test functions on any compact window.

\item \textbf{Lipschitz continuity (A2):} The functional $Q$ is Lipschitz continuous on compact windows with an explicit constant.
\end{enumerate}

These properties are summarized in Table~\ref{tab:modules}.

\begin{table}[h]
\centering
\caption{Analytic modules established in this paper.}
\label{tab:modules}
\begin{tabular}{@{}llp{7cm}@{}}
\toprule
Module & Name & Statement \\
\midrule
(T0) & Normalization & Guinand--Weil matching under $\eta = 2\pi\xi$ \\
(A1') & Density & Fej\'er--heat cone dense in $C^+_{\mathrm{even}}([-K,K])$ \\
(A2) & Continuity & $Q$ Lipschitz on compact windows \\
\bottomrule
\end{tabular}
\end{table}

We now state the main results precisely.

\begin{theorem}[T0: Guinand--Weil normalization]\label{thm:T0}
Let $\varphi_{\mathrm{GW}}$ be an even test function in the classical Guinand--Weil formulation with frequency variable $\eta$. Define $\varphi(\xi) := \varphi_{\mathrm{GW}}(2\pi\xi)$. Then the repository quadratic functional~\eqref{eq:Q-def} satisfies
\begin{equation}\label{eq:T0-match}
  Q(\varphi) = Q_{\mathrm{GW}}(\varphi_{\mathrm{GW}}),
\end{equation}
where $Q_{\mathrm{GW}}$ is the classical Guinand--Weil functional. The Archimedean densities are related by $a_*(\xi) = 2\pi \cdot a(\xi)$, accounting for the Jacobian of the coordinate change.
\end{theorem}

\begin{theorem}[A1': Fej\'er--heat cone density]\label{thm:A1}
Let $K = [-R, R]$ with $R > 0$. Define the Fej\'er--heat generators
\begin{equation}\label{eq:Phi-gen}
  \Phi_{B,t,\tau}(\xi) := \frac{1}{2}\Bigl[\Lambda_B(\xi - \tau)\, \rho_t(\xi - \tau) + \Lambda_B(\xi + \tau)\, \rho_t(\xi + \tau)\Bigr],
\end{equation}
where $\Lambda_B(x) = (1 - |x|/B)_+$ is the Fej\'er triangle of bandwidth $B$, and $\rho_t(x) = (4\pi t)^{-1/2} e^{-x^2/(4t)}$ is the heat kernel. Let $\mathcal{C}$ be the closed convex cone generated by finite non-negative combinations of $\{\Phi_{B,t,\tau}\}$ with $\tau \in [-R, R]$ and $B \geq B_0(R, \varepsilon)$. Then $\mathcal{C}$ is dense in $C^+_{\mathrm{even}}([-R,R])$ in the uniform norm.
\end{theorem}

\begin{lemma}[A2: Lipschitz continuity]\label{lem:A2}
Fix $K = [-R, R]$. For even, non-negative test functions $\Phi_1, \Phi_2$ supported in $K$, one has
\begin{equation}\label{eq:A2-Lip}
  |Q(\Phi_1) - Q(\Phi_2)| \leq L_Q(K)\, \|\Phi_1 - \Phi_2\|_\infty,
\end{equation}
where the Lipschitz constant is
\begin{equation}\label{eq:LQ}
  L_Q(K) := \|a_*\|_{L^1(K)} + \sum_{\xi_n \in K} w_Q(n) < \infty.
\end{equation}
In particular, $Q$ is continuous on $C^+_{\mathrm{even}}(K)$ with respect to the uniform topology.
\end{lemma}

\begin{remark}[Local finiteness]\label{rem:local-finite}
The finiteness of $L_Q(K)$ depends crucially on the fact that only finitely many prime nodes $\xi_n$ fall within the compact window $K$. Indeed, $\xi_n \leq R$ if and only if $n \leq e^{2\pi R}$, so the prime sum in~\eqref{eq:LQ} has at most $\lfloor e^{2\pi R} \rfloor$ terms.
\end{remark}

\subsection{Overview of the proof}

We now discuss the methods of the proof. The three modules are largely independent, but together they form the foundation for positivity arguments developed in subsequent work.

\paragraph{The normalization (T0).}
The classical Guinand--Weil functional uses the frequency variable $\eta \in \R$ with $\varphi_{\mathrm{GW}}(\pm \log n)$ appearing in the prime sum. Our coordinate system uses $\xi = \eta/(2\pi)$, so prime nodes become $\xi_n = \log n/(2\pi)$. The Jacobian $d\eta = 2\pi\, d\xi$ is absorbed into the Archimedean density, yielding $a_*(\xi) = 2\pi \cdot a(\xi)$. This normalization is consistent with the standard Fourier transform conventions
\[
  \widehat{\varphi}(\xi) = \int_\R \varphi(t)\, e^{-2\pi i t\xi}\, dt,
\]
and ensures that the weight $w_Q(n) = 2\Lambda(n)/\sqrt{n}$ at positive nodes is equivalent to placing unit weights $\Lambda(n)/\sqrt{n}$ at both $\pm\xi_n$ for even test functions.

\paragraph{The density argument (A1').}
The proof of Theorem~\ref{thm:A1} proceeds in three steps:
\begin{enumerate}
\item \emph{Mollification:} Given $f \in C^+_{\mathrm{even}}([-R,R])$, convolve with the heat kernel $\rho_t$ to obtain a smooth approximation $g = f * \rho_t$ with $\|g - f\|_\infty < \varepsilon/3$.

\item \emph{Riemann decomposition:} Approximate $g$ by a positive Riemann sum of translates of $\rho_t$, yielding $g_R(\xi) = \sum_j c_j\, \rho_t(\xi - \tau_j)$ with $c_j \geq 0$ and $\|g_R - g\|_\infty < \varepsilon/3$.

\item \emph{Fej\'er truncation:} Replace each $\rho_t(\xi - \tau_j)$ by the localized product $\Lambda_B(\xi - \tau_j)\, \rho_t(\xi - \tau_j)$. For $B$ large, the Fej\'er factor satisfies $|\Lambda_B(\xi - \tau) - 1| < \varepsilon/(3C)$ on $[-R,R]$, contributing the final $\varepsilon/3$ error.
\end{enumerate}
The symmetrization by the factor $\frac{1}{2}[\cdot + \cdot(-\cdot)]$ in~\eqref{eq:Phi-gen} ensures that the generators are even. The triangle inequality collects the three error contributions.

\paragraph{The Lipschitz bound (A2).}
The Archimedean term is controlled by
\[
  \left|\int_K a_*(\xi)\, (\Phi_1 - \Phi_2)(\xi)\, d\xi\right| \leq \|a_*\|_{L^1(K)}\, \|\Phi_1 - \Phi_2\|_\infty.
\]
For the prime term, local finiteness (Remark~\ref{rem:local-finite}) ensures
\[
  \left|\sum_{\xi_n \in K} w_Q(n)\, (\Phi_1 - \Phi_2)(\xi_n)\right| \leq \left(\sum_{\xi_n \in K} w_Q(n)\right) \|\Phi_1 - \Phi_2\|_\infty.
\]
Adding these bounds yields the Lipschitz constant $L_Q(K)$.

\begin{remark}[Physical intuition]\label{rem:physics}
The Fej\'er--heat generators can be thought of as ``frequency probes'' that sample the Weil functional at controlled resolution. The Fej\'er triangle provides bandwidth limitation (compact frequency support), while the Gaussian heat kernel provides smoothness (rapid decay in time). Together, they form an optimal family for approximation in $L^\infty$, analogous to the wavelet-like constructions used in harmonic analysis.
\end{remark}

\begin{remark}[Connection to Part~II]\label{rem:part2}
The density (A1') and continuity (A2) modules established here are combined in [Part~II] with spectral bounds on Toeplitz operators (A3) and RKHS contraction estimates to study positivity properties of $Q$. The modular structure allows each component to be developed and verified independently.
\end{remark}

\subsection{Notation}\label{sec:notation}

We collect the notation used throughout the paper.

\begin{itemize}
\item $\Lambda(n)$ denotes the von Mangoldt function: $\Lambda(p^k) = \log p$ for prime powers, zero otherwise.

\item $\xi_n = \frac{\log n}{2\pi}$ are the prime nodes for $n \geq 2$.

\item $w_Q(n) = \frac{2\Lambda(n)}{\sqrt{n}}$ are the Weil weights (evenized).

\item $w_{\mathrm{rkhs}}(n) = \frac{\Lambda(n)}{\sqrt{n}}$ are the operator weights (undoubled).

\item $w_{\max} = \sup_n w_{\mathrm{rkhs}}(n) \leq 2/e$ is the maximum operator weight.

\item $a_*(\xi) = 2\pi\bigl(\log\pi - \Re\psi(\tfrac{1}{4} + i\pi\xi)\bigr)$ is the Archimedean density.

\item $\psi(s) = \Gamma'(s)/\Gamma(s)$ is the digamma function.

\item $\Lambda_B(x) = (1 - |x|/B)_+$ is the Fej\'er triangle.

\item $\rho_t(x) = (4\pi t)^{-1/2} e^{-x^2/(4t)}$ is the heat kernel normalized to integrate to $1$.

\item $W_K = [-K, K]$ denotes the compact frequency window.

\item $C^+_{\mathrm{even}}(K)$ denotes even, non-negative continuous functions on $K$.

\item $\|\cdot\|_\infty$ denotes the uniform (supremum) norm.

\item $O(\cdot)$, $O_\varepsilon(\cdot)$, $\tilde{O}(\cdot)$ denote standard asymptotic notation.
\end{itemize}

%=============================================================================
\section{Guinand--Weil normalization (T0)}
\label{sec:T0}
%=============================================================================

In this section we establish the correspondence between our quadratic functional $Q$ and the classical Guinand--Weil formulation. The main result is that the two agree under a simple coordinate change.

\subsection{Fourier conventions}

We fix the Fourier transform convention
\begin{equation}\label{eq:FT}
  \widehat{\varphi}(\xi) = \int_\R \varphi(t)\, e^{-2\pi i t\xi}\, dt,
  \qquad
  \varphi(t) = \int_\R \widehat{\varphi}(\xi)\, e^{2\pi i t\xi}\, d\xi,
\end{equation}
using Lebesgue measure $d\xi$ on the frequency side. For even test functions, all identities reduce to cosine transforms.

\subsection{The classical Guinand--Weil functional}

Let $\varphi_{\mathrm{GW}} \in C_c(\R)$ be an even, non-negative test function on the classical frequency axis $\eta \in \R$. The \emph{Guinand--Weil functional}~\cite{Guinand1948,Weil1952} is
\begin{equation}\label{eq:QGW}
  Q_{\mathrm{GW}}(\varphi_{\mathrm{GW}}) := \int_\R \Bigl(\log\pi - \Re\psi\bigl(\tfrac{1}{4} + \tfrac{i\eta}{2}\bigr)\Bigr)\, \varphi_{\mathrm{GW}}(\eta)\, d\eta
  - \sum_{n \geq 2} \frac{\Lambda(n)}{\sqrt{n}}\, \bigl(\varphi_{\mathrm{GW}}(\log n) + \varphi_{\mathrm{GW}}(-\log n)\bigr).
\end{equation}

\subsection{Coordinate change}

On our frequency axis $\xi = \eta/(2\pi)$, define the scaled test function
\begin{equation}\label{eq:phi-scale}
  \varphi(\xi) := \varphi_{\mathrm{GW}}(2\pi\xi).
\end{equation}
The prime nodes become $\xi_n = \frac{\log n}{2\pi}$, and we define the Archimedean density
\begin{equation}\label{eq:astar}
  a(\xi) := \log\pi - \Re\psi\bigl(\tfrac{1}{4} + i\pi\xi\bigr),
  \qquad
  a_*(\xi) := 2\pi \cdot a(\xi).
\end{equation}
The factor $2\pi$ arises from the Jacobian $d\eta = 2\pi\, d\xi$.

\begin{proposition}[T0': Guinand--Weil matching]\label{prop:T0}
Under the coordinate change $\eta = 2\pi\xi$ and the scaling~\eqref{eq:phi-scale}, one has
\begin{equation}\label{eq:T0-identity}
  Q(\varphi) = Q_{\mathrm{GW}}(\varphi_{\mathrm{GW}}),
\end{equation}
where
\begin{equation}\label{eq:Q-repo}
  Q(\varphi) := \int_\R a_*(\xi)\, \varphi(\xi)\, d\xi - \sum_{n \geq 2} w_Q(n)\, \varphi(\xi_n),
  \qquad
  w_Q(n) = \frac{2\Lambda(n)}{\sqrt{n}}.
\end{equation}
\end{proposition}

\begin{proof}
We verify each term separately.

\textbf{Archimedean term.} Substituting $\eta = 2\pi\xi$ with $d\eta = 2\pi\, d\xi$:
\begin{align}
  \int_\R \Bigl(\log\pi - \Re\psi\bigl(\tfrac{1}{4} + \tfrac{i\eta}{2}\bigr)\Bigr)\, \varphi_{\mathrm{GW}}(\eta)\, d\eta
  &= \int_\R \Bigl(\log\pi - \Re\psi\bigl(\tfrac{1}{4} + i\pi\xi\bigr)\Bigr)\, \varphi(\xi)\, 2\pi\, d\xi \notag \\
  &= \int_\R a_*(\xi)\, \varphi(\xi)\, d\xi. \label{eq:arch-match}
\end{align}

\textbf{Prime term.} For even $\varphi$, we have $\varphi(\xi_n) = \varphi(-\xi_n)$, so
\[
  \varphi_{\mathrm{GW}}(\log n) + \varphi_{\mathrm{GW}}(-\log n) = \varphi(\xi_n) + \varphi(-\xi_n) = 2\varphi(\xi_n).
\]
Thus
\begin{equation}\label{eq:prime-match}
  \sum_{n \geq 2} \frac{\Lambda(n)}{\sqrt{n}}\, \bigl(\varphi_{\mathrm{GW}}(\log n) + \varphi_{\mathrm{GW}}(-\log n)\bigr)
  = \sum_{n \geq 2} \frac{2\Lambda(n)}{\sqrt{n}}\, \varphi(\xi_n)
  = \sum_{n \geq 2} w_Q(n)\, \varphi(\xi_n).
\end{equation}

Combining~\eqref{eq:arch-match} and~\eqref{eq:prime-match} yields $Q(\varphi) = Q_{\mathrm{GW}}(\varphi_{\mathrm{GW}})$.
\end{proof}

\begin{lemma}[Normalization invariance]\label{lem:norm-inv}
Different choices of Fourier normalization and node indexing yield equivalent formulations. Specifically:
\begin{enumerate}[label=(\alph*)]
\item The unitary convention~\eqref{eq:FT} versus the angular convention $\widehat{\Phi}'(\eta) = \int \Phi(x)\, e^{-i\eta x}\, dx$ induces the density rescaling $a_*(\xi) = 2\pi a(\xi)$.

\item Placing nodes at $\pm\xi_n = \pm\frac{\log n}{2\pi}$ preserves the symmetry of the sampling operator.

\item The positivity of $Q$ is independent of these technical choices.
\end{enumerate}
\end{lemma}

\begin{proof}
Each rescaling is a linear change of variable preserving the spectral structure. The node symmetry $\pm\xi_n$ is built into the Guinand--Weil formalism. The measure conversion follows from the Jacobian.
\end{proof}

\begin{remark}[Weight conventions]\label{rem:weights}
Throughout this paper, we use $w_Q(n) = 2\Lambda(n)/\sqrt{n}$ for the Weil functional (doubled weights at positive nodes for even test functions), and $w_{\mathrm{rkhs}}(n) = \Lambda(n)/\sqrt{n}$ for operator bounds in the RKHS analysis of [Part~II]. The maximum weight satisfies
\[
  w_{\max} := \sup_{n \geq 2} w_{\mathrm{rkhs}}(n) = \sup_{p \text{ prime}} \frac{\log p}{\sqrt{p}} = \frac{\log 2}{\sqrt{2}} \approx 0.490 < \frac{2}{e}.
\]
\end{remark}

%=============================================================================
\section{Fej\'er--heat cone density (A1')}
\label{sec:A1}
%=============================================================================

In this section we prove that the Fej\'er--heat generators form a dense cone in the space of even, non-negative continuous functions on compact windows.

\subsection{Fej\'er--heat generators}

\begin{definition}[Fej\'er--heat generators]\label{def:FH-gen}
For $B > 0$ (bandwidth), $t > 0$ (heat scale), and $\tau \in \R$ (center), the \emph{Fej\'er--heat atom} is
\begin{equation}\label{eq:atom}
  \Psi_{B,t,\tau}(\xi) := \Lambda_B(\xi - \tau)\, \rho_t(\xi - \tau),
\end{equation}
where $\Lambda_B(x) = (1 - |x|/B)_+$ and $\rho_t(x) = (4\pi t)^{-1/2} e^{-x^2/(4t)}$.

The \emph{symmetrized Fej\'er--heat generator} is
\begin{equation}\label{eq:gen-sym}
  \Phi_{B,t,\tau}(\xi) := \frac{1}{2}\bigl[\Psi_{B,t,\tau}(\xi) + \Psi_{B,t,\tau}(-\xi)\bigr]
  = \frac{1}{2}\bigl[\Psi_{B,t,\tau}(\xi) + \Psi_{B,t,-\tau}(\xi)\bigr],
\end{equation}
which is even by construction.
\end{definition}

\begin{remark}[Properties of generators]\label{rem:gen-props}
The Fej\'er--heat generators satisfy:
\begin{enumerate}[label=(\roman*)]
\item \emph{Non-negativity:} $\Phi_{B,t,\tau}(\xi) \geq 0$ for all $\xi \in \R$.
\item \emph{Evenness:} $\Phi_{B,t,\tau}(-\xi) = \Phi_{B,t,\tau}(\xi)$.
\item \emph{Compact support:} $\supp(\Psi_{B,t,\tau}) \subseteq [\tau - B, \tau + B]$.
\item \emph{Smoothness:} $\Phi_{B,t,\tau} \in C^\infty(\R)$ for any $t > 0$.
\item \emph{Normalization:} $\int_\R \rho_t(x)\, dx = 1$.
\end{enumerate}
\end{remark}

\begin{definition}[Fej\'er--heat cone]\label{def:FH-cone}
For a compact window $K = [-R, R]$, the \emph{Fej\'er--heat cone} is
\begin{equation}\label{eq:cone}
  \mathcal{C}_K := \overline{\Bigl\{\sum_{j=1}^n c_j\, \Phi_{B_j, t_j, \tau_j} : n \in \N,\, c_j \geq 0,\, \tau_j \in [-R, R]\Bigr\}}^{\|\cdot\|_\infty},
\end{equation}
the uniform closure of finite non-negative combinations of generators centered in $K$.
\end{definition}

\subsection{Approximation by mollification}

The first step in proving density is to approximate by convolution with the heat kernel.

\begin{lemma}[Heat kernel mollification]\label{lem:mollify}
Let $f \in C^+_{\mathrm{even}}([-R,R])$ and $\varepsilon > 0$. There exists $t_0 > 0$ such that for all $0 < t \leq t_0$, the convolution
\begin{equation}\label{eq:mollify}
  g(\xi) := \int_\R \tilde{f}(y)\, \rho_t(\xi - y)\, dy
\end{equation}
satisfies $\sup_{|\xi| \leq R} |g(\xi) - f(\xi)| < \varepsilon/3$, where $\tilde{f}$ is the extension of $f$ by zero.
\end{lemma}

\begin{proof}
Since $\rho_t$ is an approximate identity as $t \to 0^+$, the convolution $\tilde{f} * \rho_t$ converges uniformly to $\tilde{f}$ on compacts. On $[-R, R]$ we have $\tilde{f} = f$, yielding the stated bound for $t \leq t_0$ sufficiently small.
\end{proof}

\subsection{Positive Riemann sums}

\begin{lemma}[Riemann sum approximation]\label{lem:riemann}
Let $g \in C^+([-R,R])$ be non-negative and continuous, and let $\varepsilon > 0$. There exists a partition $-R = \tau_0 < \tau_1 < \cdots < \tau_N = R$ with mesh $\Delta = \max_j(\tau_{j+1} - \tau_j)$ sufficiently small such that
\begin{equation}\label{eq:riemann}
  g_R(\xi) := \sum_{j=0}^{N-1} g(\tau_j^*)\, \rho_t(\xi - \tau_j^*)\, (\tau_{j+1} - \tau_j)
\end{equation}
satisfies $\sup_{|\xi| \leq R} |g_R(\xi) - g(\xi)| < \varepsilon/3$, for some choices $\tau_j^* \in [\tau_j, \tau_{j+1}]$.
\end{lemma}

\begin{proof}
This is a standard Riemann sum approximation for the identity $g(\xi) = \int_{-R}^R \delta(\xi - \eta)\, g(\eta)\, d\eta$ regularized by $\rho_t$. The coefficients $c_j = g(\tau_j^*)(\tau_{j+1} - \tau_j) \geq 0$ since $g \geq 0$.
\end{proof}

\subsection{Fej\'er truncation}

\begin{lemma}[Fej\'er truncation error]\label{lem:fejer}
For $|\xi|, |\tau| \leq R$ and $B \geq B_0 := 6R/\varepsilon$, one has
\begin{equation}\label{eq:fejer-err}
  |\Lambda_B(\xi - \tau) - 1| \leq \frac{|\xi - \tau|}{B} \leq \frac{2R}{B} < \frac{\varepsilon}{3}.
\end{equation}
\end{lemma}

\begin{proof}
On $[-R, R]$, we have $|\xi - \tau| \leq 2R$. Since $|\xi - \tau| \leq 2R < B$, the Fej\'er function satisfies $\Lambda_B(\xi - \tau) = 1 - |\xi - \tau|/B$, and the bound follows.
\end{proof}

\subsection{Proof of Theorem~\ref{thm:A1}}

We now combine the three lemmas to prove density.

\begin{proof}[Proof of Theorem~\ref{thm:A1}]
Fix $f \in C^+_{\mathrm{even}}([-R,R])$ and $\varepsilon > 0$. We construct an approximant $h \in \mathcal{C}_K$ with $\|h - f\|_\infty < \varepsilon$.

\textbf{Step 1 (Mollification).} By Lemma~\ref{lem:mollify}, choose $t > 0$ such that $g := \tilde{f} * \rho_t$ satisfies $\|g - f\|_\infty < \varepsilon/3$ on $[-R, R]$.

\textbf{Step 2 (Riemann sums).} By Lemma~\ref{lem:riemann}, choose a partition with mesh $\Delta$ such that
\[
  g_R(\xi) = \sum_{j=0}^{N-1} c_j\, \rho_t(\xi - \tau_j^*), \qquad c_j := g(\tau_j^*)\, (\tau_{j+1} - \tau_j) \geq 0,
\]
satisfies $\|g_R - g\|_\infty < \varepsilon/3$.

\textbf{Step 3 (Fej\'er truncation).} Choose $B \geq 6R/\varepsilon$. Define the approximant
\begin{equation}\label{eq:approx-h}
  h(\xi) := \sum_{j=0}^{N-1} c_j\, \Phi_{B, t, \tau_j^*}(\xi).
\end{equation}
By Lemma~\ref{lem:fejer}, the Fej\'er factor satisfies $|\Lambda_B(\xi - \tau_j^*) - 1| < \varepsilon/(3C)$ on $[-R, R]$, where $C = \sum_j c_j = \int_K g\, d\xi$. The symmetrization in $\Phi_{B,t,\tau}$ preserves this bound.

Therefore, for $|\xi| \leq R$:
\begin{align}
  |h(\xi) - g_R^{\mathrm{sym}}(\xi)|
  &\leq \sum_j c_j\, |\Lambda_B(\xi - \tau_j^*) - 1|\, \rho_t(\xi - \tau_j^*) \notag \\
  &\leq \frac{\varepsilon}{3C} \sum_j c_j = \frac{\varepsilon}{3}. \label{eq:step3-err}
\end{align}

\textbf{Step 4 (Triangle inequality).} Since $g$ is even (as $f$ is even and $\rho_t$ is symmetric), the symmetrized Riemann sum $g_R^{\mathrm{sym}}$ satisfies $\|g_R^{\mathrm{sym}} - g\|_\infty < \varepsilon/3$. Combining:
\[
  \|h - f\|_\infty \leq \|h - g_R^{\mathrm{sym}}\|_\infty + \|g_R^{\mathrm{sym}} - g\|_\infty + \|g - f\|_\infty
  < \frac{\varepsilon}{3} + \frac{\varepsilon}{3} + \frac{\varepsilon}{3} = \varepsilon.
\]

Since $h$ is a finite non-negative combination of Fej\'er--heat generators $\Phi_{B,t,\tau_j^*}$ with $\tau_j^* \in [-R, R]$, we have $h \in \mathcal{C}_K$. Taking the uniform closure yields density.
\end{proof}

\begin{corollary}[Fixed-$t$ density]\label{cor:fixed-t}
Fix $K > 0$ and $t_0 > 0$. The cone generated by $\{\Phi_{B,t_0,\tau} : B > 0,\, |\tau| + B \leq K\}$ is dense in $C^+_{\mathrm{even}}([-K,K])$.
\end{corollary}

\begin{proof}
The proof follows the same structure, using the fixed heat scale $t_0$ throughout and choosing $B$ large depending on $\varepsilon$.
\end{proof}

\begin{remark}[Parameter dependence]\label{rem:params}
The bandwidth $B$ and heat scale $t$ required in the approximation depend on the compact $K$ and the target accuracy $\varepsilon$. In the main closure argument of [Part~III], the schedules $K \mapsto B(K)$ and $K \mapsto t(K)$ are allowed to grow with $K$; no uniform bound is needed for density.
\end{remark}

%=============================================================================
\section{Lipschitz continuity of $Q$ (A2)}
\label{sec:A2}
%=============================================================================

In this section we establish that the Weil functional $Q$ is Lipschitz continuous on compact frequency windows.

\subsection{Local finiteness of the prime sampler}

The key observation is that only finitely many prime nodes fall within any bounded interval.

\begin{lemma}[Local finiteness]\label{lem:local-finite}
Fix $K > 0$. For any even $\Phi \in C_c(\R)$ with $\supp(\Phi) \subseteq [-K, K]$, the prime sum
\begin{equation}\label{eq:prime-sum}
  \sum_{n \geq 2} w_Q(n)\, \Phi(\xi_n)
\end{equation}
involves only finitely many non-zero terms.
\end{lemma}

\begin{proof}
Since $\Phi(\xi_n) = 0$ whenever $|\xi_n| > K$, and $|\xi_n| \leq K$ if and only if $n \leq e^{2\pi K}$, only indices $n \leq \lfloor e^{2\pi K} \rfloor$ contribute. This is a finite set.
\end{proof}

\begin{remark}[Node spacing]\label{rem:spacing}
The prime nodes in $[-K, K]$ have a positive minimum spacing
\begin{equation}\label{eq:delta}
  \delta_K := \min_{\substack{m \neq n \\ \xi_m, \xi_n \in [-K,K]}} |\xi_m - \xi_n|
  \geq \frac{1}{2\pi(e^{2\pi K} + 1)},
\end{equation}
which records the lack of accumulation points. This bound is not needed for finiteness but becomes important in the RKHS analysis of [Part~II].
\end{remark}

\subsection{Lipschitz bound}

\begin{proposition}[Lipschitz continuity]\label{prop:Lip}
Let $K = [-R, R]$. For even, non-negative $\Phi_1, \Phi_2 \in C_c(K)$,
\begin{equation}\label{eq:Lip-bound}
  |Q(\Phi_1) - Q(\Phi_2)| \leq L_Q(K)\, \|\Phi_1 - \Phi_2\|_\infty,
\end{equation}
where
\begin{equation}\label{eq:LQ-def}
  L_Q(K) := \|a_*\|_{L^1(K)} + \sum_{\xi_n \in K} w_Q(n).
\end{equation}
\end{proposition}

\begin{proof}
We bound each term in $Q$ separately.

\textbf{Archimedean term.} Since $a_*$ is bounded and continuous on the compact $K$,
\begin{equation}\label{eq:arch-Lip}
  \left|\int_K a_*(\xi)\, (\Phi_1 - \Phi_2)(\xi)\, d\xi\right|
  \leq \int_K |a_*(\xi)|\, |\Phi_1(\xi) - \Phi_2(\xi)|\, d\xi
  \leq \|a_*\|_{L^1(K)}\, \|\Phi_1 - \Phi_2\|_\infty.
\end{equation}

\textbf{Prime term.} By Lemma~\ref{lem:local-finite}, only finitely many terms contribute:
\begin{equation}\label{eq:prime-Lip}
  \left|\sum_{\xi_n \in K} w_Q(n)\, (\Phi_1 - \Phi_2)(\xi_n)\right|
  \leq \sum_{\xi_n \in K} w_Q(n)\, |\Phi_1(\xi_n) - \Phi_2(\xi_n)|
  \leq \Bigl(\sum_{\xi_n \in K} w_Q(n)\Bigr)\, \|\Phi_1 - \Phi_2\|_\infty.
\end{equation}

Adding~\eqref{eq:arch-Lip} and~\eqref{eq:prime-Lip} yields the result.
\end{proof}

\begin{corollary}[Continuity of $Q$]\label{cor:cont}
The Weil functional $Q : C^+_{\mathrm{even}}(K) \to \R$ is continuous with respect to the uniform topology.
\end{corollary}

\begin{proof}
Lipschitz continuity implies continuity.
\end{proof}

\subsection{Tail control for Fej\'er--heat windows}

When Fej\'er--heat windows have slight leakage outside the compact $K$, the Gaussian decay provides explicit tail bounds.

\begin{lemma}[Gaussian tail bound]\label{lem:tail}
Let $\Phi$ be a Fej\'er--heat window with heat scale $t > 0$. For any $N \geq e^2$, the tail contribution satisfies
\begin{equation}\label{eq:tail}
  \mathrm{Tail}(t; N) := \sum_{\substack{\xi_n \notin K \\ n > N}} w_Q(n)\, \Phi(\xi_n)
  \leq \frac{e^{-t(\log N)^2}}{t}.
\end{equation}
\end{lemma}

\begin{proof}
For Fej\'er--heat windows, $\Phi(\xi) \leq e^{-4\pi^2 t \xi^2}$. At nodes $\xi_n = \frac{\log n}{2\pi}$,
\[
  \Phi(\xi_n) \leq e^{-t(\log n)^2}.
\]
The sum over $n > N$ is bounded by the integral
\[
  \sum_{n > N} \frac{\log n}{\sqrt{n}}\, e^{-t(\log n)^2}
  \leq \int_{\log N}^\infty y\, e^{-ty^2}\, dy
  = \frac{1}{2t}\, e^{-t(\log N)^2}.
\]
Accounting for the factor $2$ in $w_Q(n)$ yields the stated bound.
\end{proof}

\begin{remark}[Leakage control]\label{rem:leakage}
The tail bound~\eqref{eq:tail} shows that Gaussian leakage outside $K$ is exponentially small in $t(\log N)^2$. For the main positivity argument in [Part~III], this leakage is absorbed into the continuity budget of the A2 module.
\end{remark}

%=============================================================================
\section{Discussion}
\label{sec:discussion}
%=============================================================================

\subsection{Summary of results}

We have established three foundational properties of the Weil quadratic functional:

\begin{enumerate}
\item \textbf{Normalization (T0):} The correspondence $\eta = 2\pi\xi$ matches our functional $Q$ with the classical Guinand--Weil formulation, with the Jacobian absorbed into the Archimedean density $a_*(\xi) = 2\pi a(\xi)$.

\item \textbf{Density (A1'):} The Fej\'er--heat cone---finite non-negative combinations of products of Fej\'er triangles with Gaussian heat kernels---is dense in $C^+_{\mathrm{even}}([-K,K])$ for any compact window $K$.

\item \textbf{Continuity (A2):} The functional $Q$ is Lipschitz continuous on compact windows with explicit constant $L_Q(K) = \|a_*\|_{L^1(K)} + \sum_{\xi_n \in K} w_Q(n)$.
\end{enumerate}

These properties form the analytic foundation for studying positivity of $Q$.

\subsection{Connection to Part~II}

The results established here are prerequisites for the spectral analysis in [Part~II], where:

\begin{itemize}
\item The \textbf{Toeplitz barrier (A3)} provides a positive lower bound $c_* = 11/10$ on the Archimedean symbol $P_A(\theta)$ for Fej\'er--heat windows with bandwidth $B \geq 3$.

\item The \textbf{RKHS prime cap} controls the operator norm of the prime sampling operator $T_P$ in the heat kernel RKHS, yielding $\|T_P\| \leq \rho(t) < c_*/4$ for $t \geq t_*$.
\end{itemize}

The combination of these bounds yields the spectral gap
\[
  \lambda_{\min}(T_M[P_A] - T_P) \geq \frac{c_*}{4} > 0,
\]
from which positivity $Q(\Phi) \geq 0$ follows for Fej\'er--heat windows. The density (A1') and continuity (A2) modules then extend this positivity from the generators to all even, non-negative test functions.

\subsection{Connection to Part~III}

The full synthesis of modules (T0), (A1'), (A2), (A3), and (RKHS) to establish positivity of $Q$ on the Weil cone $\mathcal{W} = \bigcup_K \mathcal{W}_K$ appears in [Part~III]. The compact-by-compact approach---establishing $Q \geq 0$ on each $W_K$ and then taking the union---requires all the machinery developed here and in [Part~II].

\subsection{Open questions}

Several natural questions arise from this work:

\begin{enumerate}
\item \textbf{Optimal constants:} Can the Lipschitz constant $L_Q(K)$ be sharpened? The current bound may not be tight.

\item \textbf{Alternative generators:} Are there other families of generators (e.g., wavelets, prolate spheroidal functions) that achieve density with better approximation rates?

\item \textbf{Quantitative density:} Can one obtain explicit error bounds $\|h - f\|_\infty \leq \varepsilon$ in terms of the parameters $(B, t, N)$ of the approximant?
\end{enumerate}

These questions may be relevant for optimizing the numerical aspects of positivity verification.

%=============================================================================
% Bibliography
%=============================================================================

\begin{thebibliography}{99}

\bibitem{BoettcherSilbermann2006}
A.~B\"ottcher and B.~Silbermann,
\emph{Analysis of Toeplitz Operators},
2nd ed., Springer Monographs in Mathematics, Springer, Berlin, 2006.

\bibitem{GrenanderSzego1958}
U.~Grenander and G.~Szeg\H{o},
\emph{Toeplitz Forms and Their Applications},
University of California Press, Berkeley, 1958.

\bibitem{Guinand1948}
A.~P.~Guinand,
``A summation formula in the theory of prime numbers,''
\emph{Proc.\ London Math.\ Soc.}\ \textbf{50} (1948), 107--119.

\bibitem{NISTDLMF}
NIST Digital Library of Mathematical Functions,
\url{https://dlmf.nist.gov/}, Release 1.1.0.

\bibitem{SteinShakarchi2003}
E.~M.~Stein and R.~Shakarchi,
\emph{Fourier Analysis: An Introduction},
Princeton University Press, Princeton, 2003.

\bibitem{Szego1952}
G.~Szeg\H{o},
``On certain Hermitian forms associated with the Fourier series of a positive function,''
\emph{Comm.\ S\'em.\ Math.\ Univ.\ Lund, Tome Suppl\'ementaire} (1952), 228--238.

\bibitem{Titchmarsh1986}
E.~C.~Titchmarsh,
\emph{The Theory of the Riemann Zeta-Function},
2nd ed.\ (revised by D.~R.~Heath-Brown), Oxford University Press, Oxford, 1986.

\bibitem{Weil1952}
A.~Weil,
``Sur les 'formules explicites' de la th\'eorie des nombres premiers,''
\emph{Comm.\ S\'em.\ Math.\ Univ.\ Lund, Tome Suppl\'ementaire} (1952), 252--265.

\end{thebibliography}

\end{document}
